\documentclass[a4paper,10pt]{article}
\usepackage[english]{babel}
\usepackage{url}

%\usepackage[utf8]{inputenc}
\usepackage[T1]{fontenc}
\usepackage{babel}
%\usepackage{graphicx}
\usepackage{xspace}
\usepackage{color}
\usepackage{setspace}
\usepackage{geometry}
%\usepackage{helvet}
\usepackage{calc}
\usepackage{ifthen}


\pagestyle{empty}

%%%%%
%%%%%

% Package pour la gestion des couleurs
%
\usepackage{color}

% Les commandes pour passer de la version nb à la version couleur 
% (et réciproquement)
%
\newboolean{enNoir}
\setboolean{enNoir}{true}
\newboolean{enCouleur}
\setboolean{enCouleur}{false}
\newcommand{\enNoir}{\definecolor{couleur}{rgb}{0,0,0}
  \setboolean{enCouleur}{false}
  \setboolean{enNoir}{true}}
\newcommand{\enGris}{\definecolor{couleur}{rgb}{0,0,0}
  \setboolean{enCouleur}{false}
  \setboolean{enNoir}{false}}
% Définition du bleu ULP
\newcommand{\enCouleur}{\definecolor{couleur}{rgb}{0,0.2,0.6}
  \setboolean{enCouleur}{true}
  \setboolean{enNoir}{false}}
% La version par défaut en en noir
%
\enNoir

%%%%%%
% Gestion de pdf
%%%%%%
%--------------------------------------
%pour compiler avec latex ou pdflatex
\newif\ifpdf
\ifx\pdfoutput\undefined
  \pdffalse
\else
  \pdfoutput=1
  \pdftrue
\fi
%--------------------------------------
%pour eviter de generer des fonts de type 3
\ifpdf 
\usepackage{pslatex}
\usepackage[pdftex]{graphicx}
%\DeclareGraphicsExtensions{.gif}
\else
\usepackage[dvips]{graphicx}
\fi
%--------------------------------------


\addtolength{\parskip}{\baselineskip}

\newlength{\firstcolumnwidth}
\setlength{\firstcolumnwidth}{1.65in}
\newlength{\inbetween}
\setlength{\inbetween}{5mm}
\newlength{\headermargin}
\setlength{\headermargin}{49mm}

\geometry{left=0.47in+\firstcolumnwidth+\inbetween,right=1in,top=15mm,bottom=20mm,noheadfoot}

\newenvironment{letter}[3][]
{
\thispagestyle{empty}

\newsavebox{\dimargin}
\sbox{\dimargin}{%
  \begin{minipage}[t]{\firstcolumnwidth}
    \vspace{0pt}
    \sffamily
    \begin{flushright}
      %%% Entête du DI
      \newsavebox{\dihead}
      \sbox{\dihead}{%
        \begin{minipage}{\linewidth}
          \flushright
          {\bfseries
            \color{couleur}
            \begin{spacing}{1}
              \large{Laboratoire}
              
              \large{des sciences de l'ingénieur, }
            
              \large{de l'informatique et} 

              \large{de l'imagerie (Icube)} 
        
              \normalsize{UMR 7537} 
            \end{spacing}
            }
          \vspace{-2mm}
          
          \rule{36.5mm}{0.5pt}
          
%          \vspace{-4mm}
          %%% Inclusion du Logo ULP
          \ifthenelse{\boolean{enNoir}}
          {
            \flushright{\includegraphics[height=19mm]{logo-uds-couleur-800x433}}
          }
          {
            \ifthenelse{\boolean{enCouleur}}
            {
              \includegraphics[height=19mm]{}
              }
            {
              \includegraphics[height=19mm]{logo-uds-couleur-800x433}
              }
            }
          %%% Fin inclusion logo
        \end{minipage}%
        }
      \newlength{\diheadheight}
      \newlength{\diheaddepth}
      \newlength{\diheadspace}
      \settoheight{\diheadheight}{\usebox{\dihead}}
      \settodepth{\diheaddepth}{\usebox{\dihead}}
      \setlength{\diheadspace}{\diheadheight+\diheaddepth}
      %
      %
      %%% Coordonnées de l'expéditeur
      \newsavebox{\sender}
      \sbox{\sender}{%
        \begin{minipage}{\linewidth}
          \flushright
          \begin{spacing}{1.5}
            
            \textbf{\small{\fromname}}

            \textbf{\footnotesize{\fonction}}
            
            \scriptsize{Tel. : \phoneespla}
            
            \scriptsize{\phoneillkirch}
            
            \scriptsize{\emailadd}
          \end{spacing}
        \end{minipage}%
        }
      \newlength{\senderheight}
      \newlength{\senderdepth}
      \newlength{\senderspace}
      \settoheight{\senderheight}{\usebox{\sender}}
      \settodepth{\senderdepth}{\usebox{\sender}}
      \setlength{\senderspace}{\senderheight+\senderdepth}
      %
      %
      %%% Coordonnées du DI
      \newsavebox{\difoot}
      \sbox{\difoot}{%
        \begin{minipage}{\linewidth}
          \flushright
          \begin{spacing}{1.25}
            \footnotesize
            
            \textbf{Laboratoire Icube}
            
            
            {Pôle API -Bd S. Brant}
            
            {F-67400 Illkirch}
            
            {Tél. : +33 (0) 3 90 24 45 42}
            
            {Fax : +33 (0) 3 90 24 45 47}
            
            {http\string://icube.unistra.fr}
          \end{spacing}
        \end{minipage}%
        }
      \newlength{\difootheight}
      \newlength{\difootdepth}
      \newlength{\difootspace}
      \settoheight{\difootheight}{\usebox{\difoot}}
      \settodepth{\difootdepth}{\usebox{\difoot}}
      \setlength{\difootspace}{\difootheight+\difootdepth}
      %
      %
      \newlength{\devantsender}
      \setlength{\devantsender}{57mm}
%
      \newlength{\devantdifoot}
%
      \setlength{\devantdifoot}{\textheight-\diheadspace-\devantsender-\senderspace-\difootspace-8mm-13pt}
%
      \vspace{8mm}

      \usebox{\dihead}

      \vspace{\devantsender}

      \usebox{\sender}

      \vspace{\devantdifoot}

      \usebox{\difoot}
    \end{flushright}
  \end{minipage}%
}

\begin{spacing}{1}

  \noindent
  \begin{picture}(0,0)
    \setlength{\unitlength}{\firstcolumnwidth+\inbetween}
    \put(-1,0){\usebox{\dimargin}}
  \end{picture}%
%
\begin{minipage}[t]{\headermargin}

\vspace{26mm}

~
\end{minipage}%
\begin{minipage}[t]{\linewidth-\headermargin}
  %Strasbourg, le 30 octobre 2012
  Strasbourg, le \today
\end{minipage}

\noindent
\begin{minipage}[t]{\headermargin}
\vspace{35mm}
~
\end{minipage}%
%-----------------------------------------
% Champ destinataire (#2=To/A , #3 <nom>)   
%-----------------------------------------
\ifthenelse{\equal{#3}{}}{%
\begin{minipage}[t]{\linewidth-\headermargin}

\vspace{2mm}

\end{minipage}
}%
{%
\begin{minipage}[t]{\linewidth-\headermargin} #2
  
  

  \vspace{2mm}

  #3
}
\end{minipage}
\ifthenelse{\equal{#1}{}} {} {Objet : {#1}}
\vspace{20mm}
}
{
%\vspace{20mm}

\newlength{\taillesignature}
\settowidth{\taillesignature}{\fromname}
\newlength{\possignature}
\setlength{\possignature}{0.667\linewidth-0.5\taillesignature}

\hspace{\possignature}\fromname


\end{spacing}
}



\newcommand*{\name}[1]{\def\fromname{#1}}
\name{}
\newcommand*{\fonction}[1]{\def\fonction{#1}}
\newcommand*{\telephoneesplanade}[1]{\def\phoneespla{#1}}
\telephoneesplanade{}
\newcommand*{\telephoneillkirch}[1]{\def\phoneillkirch{#1}}
\telephoneillkirch{}
\newcommand*{\email}[1]{\def\emailadd{#1}}
\email{}
\newcommand*{\objet}[1]{\def\emailadd{#1}}




../letter-ensiie.sty

\begin{document}


\name{Prof. Dr. Stéphane \textsc{Genaud}}
\fonction{Professor in Computer Science}
\telephoneillkirch{+33 (0)3 69 200 231}
\email{genaud@ensiie.fr}

\begin{letter}[Recommendation of Ming Mao]%
{to}{the USCIS}
 
To the USCIS:

Cloud computing is catalyzing a series of changes with regard to the consumption
of information technology.  It allows  for shared resources, including networks,
servers   and   applications,  which   signifies   greatly   reduced  costs   of
operation. Already,  business and  economic practices have  shifted in  order to
incorporate cloud  technology and all its  benefits. These shifts are  part of a
market   in    “hypergrowth”   expected    to   reach    a   value    of   about
\$190   billion   by  2020,   according   to   the  U.S.    International   Trade
Administration. However, cloud computing  currently faces some limitations, many
of which Dr. Ming Mao has addressed and successfully overcome. His contributions
are particularly  relevant for  U.S.  businesses,  who had  spent more  than \$13
billion on  cloud computing solutions  at the  beginning of 2014.   Further, the
great  majority  of  global  firms   developing  and  offering  cloud  computing
strategies are U.S.-based. Dr. Ming Mao’s  continued research in the field is of
great importance  precisely for these  reasons. As the  very first to  develop a
standard  evaluative framework  for  virtual machine  (VM)  performance, he  has
contributed  advances of  great  importance for  cloud  computing processes  and
established himself as a researcher at the top of the field.  

Please  allow me  to introduce  myself.   I have  been a  Professor in  Computer
Science  at  ENSIIE  (National  School  of Computer  Science  for  Industry  and
Business) Engineering School in Computer Science since 2014, and Director of the
Strasbourg campus since spring 2013. Previously, I was an Assistant Professor in
Computer Science at the University of  Strasbourg. My research has a heavy focus
on  applications of  parallel and  distributed computing  to clusters,  grids or
clouds. I  am an  affiliate of  the Icube research  laboratory ICPS  group, with
which  I conduct  research on  cloud and  grid computing.  I know  Dr. Mao  only
through close  readings of  his publications,  which have  been included  in the
proceedings of high-profile IEEE and  ACM conferences. For a computer scientist,
publication at  these venues is a  hallmark of distinction as  many accept fewer
than 20\%  of submitted  papers. Some of  the conferences at  which Dr.  Mao has
published research include  the International Conference on  Grid Computing, the
International  Conference  on  Cloud  Computing (CLOUD)  and  the  International
Symposium on Parallel and Distributed Processing.   

It was at a CLOUD conference that  Dr. Mao published the results of his thorough
study of  VM performance.  He studied  VM startup and  shutdown for  three cloud
providers, Amazon  EC2, Windows  Azure and Rackspace,  according to  all factors
involved  in VM  performance. Dr.  Mao rigorously  defined metrics  according to
which he  measured and  compared the  performance of  the three  providers.  His
experiments revealed  that Rackspace  has the  greatest rate  of failure  of the
three providers, 8\%.  By contrast, Amazon EC2  has a failure rate  of 0.8\% and
Windows Azure has a failure rate of  0.4\%. He also established that the size of
the operating  system image  affects the  VM startup time  in a  positive linear
fashion for  all three  cloud providers.  He found that  the instance  type also
affects the VM startup time according  to a positive linear function for Windows
Azure and  for Rackspace, but  the two variables  are not correlated  for Amazon
EC2.  For all three  platforms, VM startup times are independent  of the time of
day and of the location of data centers except for a single data center used for
Amazon EC2. In this particular case, Dr.  Mao observed startup times around 20\%
longer than  the rest.  When recruiting  a pool of  VMs, Dr.  Mao found  that VM
startup time is fairly constant across all  instances for Amazon EC2 but not for
Windows Azure,  where the  last instance  sometimes requires  significantly more
time than the first.  He also established that the VM release time is unaffected
by the  instance type,  by the  OS image  size or  by the  location of  the data
center. Through  his studies of  these fundamental computing resources,  Dr. Mao
has  greatly   benefitted  subsequent   study  of   VM  performance   and  cloud
configuration.  He has  contributed to progressively more  efficient and elastic
cloud usage plans.  

Dr. Mao’s VM performance study has been  cited more than 150 times.  I published
an article  in the  27th IEEE International  Parallel \&  Distributed Processing
Symposium in which I describe Dr. Mao’s discovery of the relationship between VM
startup time and OS image size and detail the importance of this discovery.  The
implementation of a preemptive mechanism  to mitigate this phenomenon allows for
faster cloud computation and online application functions. Many others have used
Dr.  Mao’s  results to  configure more advanced  cloud applications  and conduct
experiments  and  simulations.   Notably,   researchers  from  Argonne  National
Laboratory directly  applied Dr.   Mao’s cloud environment  characterizations to
conduct simulations of VM overhead. Others  at Fermi National lab also cited Dr.
Mao’s work with virtual machine overhead as one of the few to note its impact on
resource allocation  during cloud  computation.  Dr.  Mao’s  work has  also been
implemented  and  confirmed by  researchers  at  the  University of  Bamberg  in
Germany, at the National Institute of  Technology in India, at the University of
Melbourne in Australia and at Tsinghua University  in China. These are but a few
prominent examples  of more  than 660  total citations  Dr. Mao’s  articles have
received. This  figure serves as  a final,  unambiguous indicator of  Dr.  Mao’s
position at the top of the international field.

Yours congenially,




\end{letter}
\begin{flushright}
\includegraphics[width=.30\textwidth]{signgenaud.jpg}
\end{flushright}
%\newpage
%\bibliographystyle{plain}
%\bibliography{julien_blaise_ED}


\end{document}

\end{letter}
\begin{flushright}
\includegraphics[width=.20\textwidth]{signgenaud.jpg}
\end{flushright}
%\newpage
%\bibliographystyle{plain}
%\bibliography{julien_blaise_ED}


\end{document}
